\subsection{LLMs and SLMs for traffic signal control}\label{sec:bg-llmtsc}

LLMLight/LightGPT \cite{lai2023llmlight}, the closest prior work, distils GPT-4 rationales into LoRA-tuned students down to 0.5B, yet serves them on four A800-80GB GPUs and evaluates single-run on CityFlow, without seeds or confidence intervals. Traffic-R1 \cite{wang2025trafficr1} trains a 3B model with GRPO, with a ten-intersection production deployment and evidence that SFT degrades the base model's general benchmarks; its own future work concedes edge deployment remains unrealised, the artefact this report builds. DGLight \cite{dglight2026} (a dissertation preprint, not peer-reviewed) tunes an 8B model on dual H100s and calibrates the field: the LLM-versus-MaxPressure travel-time gap it and its peers report spans 1.1--4.2\%, narrower than the seed variance I measure on the corridor topologies and that no competitor tests for. CoLLMLight \cite{collmlight2025} scales to 196 intersections on server GPUs and finds its fine-tuned 8B model beating a stock 671B model, the single data point for distillation over scale that motivated my central bet. CuraLight \cite{curalight2026} is the one fine-tuning work on SUMO, but its 12B model needs roughly 44GB of VRAM against my 6GB budget. iLLM-TSC \cite{pang2024illmtsc} inverts my safety topology on a single synthetic four-way intersection, an RL agent proposing and a cloud GPT-4 vetoing, with a retry loop falling back \emph{silently} to the unvetted action; its degraded-communications threat model is a scenario axis I do not test. LATS \cite{lats2026} deploys nothing generative, distilling an LLM teacher into MLP/GRU students, and its pure-LLM baselines perform worst, corroborating my finding that stock small models are weak deciders. Across all seven, every system serves from cloud or server GPUs, every headline table is single-run or at best $n\le10$ without statistical tests, none pairs control with an audit mechanism, and none evaluates on measured UK topologies.
